\section{Metodologia}
\label{sec:metodologia}

A pesquisa é aplicada, de abordagem quantitativa, e segue as seis atividades da DSR \cite{peffers2007design}: \textbf{identificação do problema e motivação} (Seção~\ref{sec:introducao}); \textbf{objetivos da solução}, isto é, quantificar disparidades, explicar suas origens (SHAP) e operar em infraestrutura restrita; \textbf{projeto e desenvolvimento} do pipeline em lote e da interface desacoplada (Seção~\ref{sec:artefato}); \textbf{demonstração}, pela aplicação às edições de 2019 a 2021 do ENEM; \textbf{avaliação}, por injeção sintética de viés com controle negativo (Seção~\ref{sec:injecao}); e \textbf{comunicação}, por este artigo, pelo código publicado e pelo painel.

\subsection{Dados e Preparação}
\label{sec:dados}

O artefato usa os microdados públicos e anonimizados das edições de \textbf{2019, 2020 e 2021} do ENEM, publicados pelo INEP. O recorte abrange a última edição anterior à pandemia de COVID-19, a edição aplicada em meio à suspensão das aulas presenciais (com provas em janeiro e fevereiro de 2021) e a edição subsequente ao período mais agudo do ensino remoto. Do arquivo de participantes, que ocupa entre 1,5 e 2,4~GB por edição, são extraídas apenas as colunas de presença e de notas, os atributos protegidos raça/cor (\texttt{TP\_COR\_RACA}: não declarada, branca, preta, parda, amarela e indígena) e renda familiar mensal (\texttt{Q006}, faixas A, nenhuma renda, a Q, acima de 20 salários mínimos), e os preditores escolaridade do pai e da mãe (\texttt{Q001}, \texttt{Q002}), computador na residência (\texttt{Q024}) e acesso à internet (\texttt{Q025}). A leitura é feita em blocos de 100.000 linhas, para respeitar o limite de 8~GB de memória, e o resultado é persistido no formato colunar Apache Parquet.

São mantidos apenas os participantes presentes nas quatro provas objetivas, pois a ausência de nota por não comparecimento não expressa desempenho; notas zero na redação de participantes presentes refletem anulação ou desempenho e são mantidas. Restam 3.701.910 participantes em 2019, 2.588.681 em 2020 e 2.238.107 em 2021. Valores ausentes no questionário são imputados pela moda de cada coluna na população filtrada da edição. A edição de 2019 não tem ausentes, e a de 2021 tem um único registro com respostas ausentes; em 2020, 27.377 participantes (1,06\%) não responderam a nenhuma das cinco questões usadas. Como a imputação precede a partição, o teste participa do cálculo da moda, mas o efeito é desprezível: o teste representa entre 0,27\% e 0,45\% da população, e a vantagem da moda sobre a segunda categoria supera 10.000 registros em todas as colunas e edições. Em 2020, porém, a ausência provavelmente não é aleatória, e atribuir a moda cria um perfil socioeconômico presumido.

O desempenho é sintetizado pela Nota Média Geral (NMG), média das notas de Ciências da Natureza, Ciências Humanas, Linguagens e Códigos, Matemática e Redação, dicotomizada pela mediana $\tilde{m}_a$ da população filtrada de cada edição $a$, calculada antes da amostragem:
\begin{equation}
\mathrm{NMG} = \frac{N_{\mathrm{CN}} + N_{\mathrm{CH}} + N_{\mathrm{LC}} + N_{\mathrm{MT}} + N_{\mathrm{RED}}}{5}, \qquad
Y = \begin{cases} 1, & \text{se } \mathrm{NMG} \geq \tilde{m}_a,\\ 0, & \text{caso contrário.} \end{cases}
\label{eq:alvo}
\end{equation}
O alvo é \textbf{relativo a cada edição}: $Y = 1$ indica desempenho igual ou superior ao do participante mediano daquele ano, e não aprovação, ingresso no ensino superior ou um patamar absoluto de proficiência. A mediana produz classes balanceadas, dispensa um limiar normativo de ``bom desempenho'', inexistente para a NMG, e permite aplicar diretamente as métricas de equidade baseadas em taxas de seleção e de erro. Por fim, extrai-se de cada edição uma amostra de \textbf{50.000 participantes}, estratificada por \texttt{TP\_COR\_RACA} com semente 42, o que preserva a proporção dos grupos minoritários.

\subsection{Modelo-Alvo da Auditoria}
\label{sec:modelo}

O objeto da auditoria é um classificador XGBoost \cite{chen2016xgboost} que recebe as seis variáveis socioeconômicas, codificadas por \textit{One-Hot Encoding} em 46 variáveis binárias. A inclusão deliberada de raça/cor como preditor simula um modelo construído sem salvaguardas de equidade, cenário em que a auditoria é mais necessária. Cada amostra é particionada, de forma estratificada pelo alvo e com semente 42, em \textbf{80\% para treino (40.000) e 20\% para teste (10.000)}. Os hiperparâmetros são selecionados por \textit{Grid Search} com validação cruzada ($k=5$) apenas no treino, com a AUC-ROC como critério, sobre a grade: número de árvores $\{50, 100\}$, profundidade máxima $\{3, 5\}$ e taxa de aprendizado $\{0{,}05;\ 0{,}1\}$. A grade é deliberadamente compacta: busca-se um modelo representativo da prática corrente, e não o de maior desempenho possível, e uma grade extensa multiplicaria o custo dos retreinamentos do experimento de injeção.

\subsection{Protocolo de Auditoria}
\label{sec:metricas}

Toda a auditoria é calculada \textbf{exclusivamente sobre o conjunto de teste} de cada edição ($n = 10.000$), que o motor de auditoria reconstrói com a mesma partição e semente do treino; a reconstrução é verificada pela coincidência exata das métricas de desempenho. O isolamento evita que exemplos memorizados pelo modelo subestimem as disparidades.

Seja $\hat{Y}$ a predição, $Y$ o rótulo real e $G$ um atributo protegido com subgrupos $g$. A taxa de seleção é $\mathrm{TS}_g = P(\hat{Y} = 1 \mid G = g)$, e a Paridade Demográfica exige $\mathrm{TS}_g = \mathrm{TS}_{g'}$ para todo par de subgrupos. O grau de violação é medido pelo Impacto Disparate \cite{feldman2015certifying}:
\begin{equation}
\mathrm{DI}_g = \frac{\mathrm{TS}_g}{\mathrm{TS}_r}, \qquad r = \arg\max_{g}\ \mathrm{TS}_g.
\label{eq:di}
\end{equation}
O grupo de referência $r$ é determinado empiricamente como o de maior taxa de seleção, o que dispensa julgamentos prévios sobre qual grupo é privilegiado (relevante para as 17 faixas de renda) e mantém $\mathrm{DI}_g \in [0, 1]$. Valores $\mathrm{DI}_g < 0{,}8$ disparam um alerta segundo a regra dos quatro quintos \cite{eeoc1978uniform}, empregada como \textbf{indicador de triagem de disparidade estatística}, e não como teste jurídico de discriminação: o alerta sinaliza que a disparidade merece investigação, mas não estabelece, por si só, ilicitude ou intenção discriminatória no ordenamento brasileiro.

Para distinguir o viés já presente nos dados daquele produzido pelo modelo, calcula-se o DI dos rótulos reais e a amplificação:
\begin{equation}
\mathrm{DI}_{y,g} = \frac{P(Y=1 \mid G=g)}{P(Y=1 \mid G=r)}, \qquad \Delta_g = \mathrm{DI}_{\hat{y},g} - \mathrm{DI}_{y,g},
\label{eq:amplificacao}
\end{equation}
em que $\mathrm{DI}_{\hat{y},g} \equiv \mathrm{DI}_g$; valores $\Delta_g < 0$ indicam que o classificador \textbf{amplifica} a disparidade dos dados (nas tabelas, o índice do grupo é omitido; $\Delta$ antes do nome de uma métrica, como em $\Delta\mathrm{AUC}$, indica a variação em relação à linha de base). Para a Igualdade de Oportunidades e as chances equalizadas \cite{hardt2016equality}, reportam-se, por subgrupo,
\begin{equation}
\mathrm{TPR}_g = P(\hat{Y} = 1 \mid Y = 1, G = g), \qquad \mathrm{FPR}_g = P(\hat{Y} = 1 \mid Y = 0, G = g),
\label{eq:tpr}
\end{equation}
comparadas às do mesmo grupo de referência $r$. No ENEM, uma $\mathrm{TPR}_g$ baixa significa que poucos candidatos do grupo $g$ com NMG igual ou superior à mediana são reconhecidos pelo modelo.

Cada estimativa é acompanhada de um intervalo de confiança de 95\% (IC95\%) por \textit{bootstrap} não paramétrico \cite{efron1993bootstrap}: o teste é reamostrado com reposição $B = 2.000$ vezes, com semente 42, e em cada réplica são recalculados DI, TPR, FPR, $\mathrm{DI}_{y,g}$ e $\Delta_g$, mantendo fixa a identidade do grupo de referência. Os limites são os percentis 2,5 e 97,5, e registra-se a proporção de réplicas com $\mathrm{DI}_g < 0{,}8$, que indica a estabilidade do alerta. Como o modelo permanece fixo, os intervalos capturam a variabilidade amostral do teste, mas não a decorrente do treinamento.

Os valores SHAP (Equação~\ref{eq:shapley}) são calculados pelo TreeSHAP nativo em C++ da biblioteca XGBoost, em escala log-odds. A importância global de cada atributo é $I_i = \frac{1}{n}\sum_{j=1}^{n} |\phi_i(f, x_j)|$ sobre as $n$ instâncias de teste, e a explicação local decompõe a predição de um candidato em gráficos de cascata (\textit{waterfall plots}).

\subsection{Validação por Injeção Sintética de Viés}
\label{sec:injecao}

Em dados reais, a magnitude verdadeira do viés é desconhecida. Para avaliar a \textbf{sensibilidade} da ferramenta, um viés de magnitude conhecida é inserido nos rótulos de treino, simulando discriminação histórica nos dados \cite{wick2019unlocking}. Seja $\mathcal{P}_{g^{*}}$ o conjunto de rótulos positivos de treino do grupo-alvo $g^{*}$. Para uma intensidade $p$, sorteia-se sem reposição um subconjunto $\mathcal{F} \subseteq \mathcal{P}_{g^{*}}$ com $|\mathcal{F}| = \operatorname{round}(p \cdot |\mathcal{P}_{g^{*}}|)$, cujos rótulos são rebaixados de 1 para 0. Como rebaixar rótulos reduz, por si só, a prevalência global de positivos, cada cenário de injeção é pareado com um \textbf{controle negativo}, em que o mesmo número $|\mathcal{F}|$ de rótulos é rebaixado por sorteio entre todos os positivos do treino. O efeito líquido atribuível ao direcionamento é
\begin{equation}
\mathrm{EL}_{g^{*}}(p) = \left[\mathrm{DI}^{\text{inj}}_{g^{*}}(p) - \mathrm{DI}_{g^{*}}(0)\right] - \left[\mathrm{DI}^{\text{ctrl}}_{g^{*}}(p) - \mathrm{DI}_{g^{*}}(0)\right].
\label{eq:efeito_liquido}
\end{equation}

Os grupos-alvo são Parda, o maior grupo não branco, e Indígena, o menor e o de menor taxa de seleção em 2019 e 2020, com $p \in \{0\%, 10\%, 25\%, 50\%\}$ e cinco repetições por combinação, com geradores derivados da semente 42: $3 \times 2 \times 4 \times 2 \times 5 = 240$ execuções. O modelo é reajustado com os \textbf{mesmos hiperparâmetros} do modelo de produção, sem nova busca em grade, e avaliado no conjunto de teste \textbf{original}, com o mesmo motor de equidade da auditoria. São testadas três hipóteses:
\begin{description}
    \item[H1 (sensibilidade):] no cenário de injeção, o DI médio do grupo-alvo decresce monotonicamente com $p$;
    \item[H2 (especificidade):] o efeito líquido é negativo, $\mathrm{EL}_{g^{*}}(p) < 0$, para toda intensidade $p > 0$, em cada grupo e edição;
    \item[H3 (baixa sensibilidade das métricas convencionais):] a AUC-ROC varia pouco entre as intensidades, quando comparada à variação do DI.
\end{description}

O trabalho usa exclusivamente microdados públicos e anonimizados, sem identificadores diretos. O modelo treinado tem finalidade exclusiva de objeto de auditoria e não deve ser empregado em decisões sobre candidatos reais.
